% eksperymentalne tłumaczenie na włoski

%

\section{Punkty szczególne} % lang-pl
\section{Punti notevoli} % lang-it

Z każdym trójkątem związane są pewne specjalne punkty, internetowa lista \emph{Encyclopedia of Triangle Centers} wymieni ich co najmniej 68 547. % lang-pl
My nie mamy aż tyle miejsca, więc ograniczymy się do najważniejszych: % lang-pl
A ogni triangolo sono associati alcuni punti notevoli; l'elenco online \emph{Encyclopedia of Triangle Centers} ne riporta almeno 68\,547. % lang-it
Non abbiamo però così tanto spazio, quindi ci limiteremo ai più importanti: % lang-it
\begin{itemize}
\item $X(1)$ to środek okręgu wpisanego, % lang-pl
\item $X(2)$ to środek ciężkości, % lang-pl
\item $X(3)$ to środek okręgu opisanego, % lang-pl
\item $X(4)$ to ortocentrum. % lang-pl
\item $X(1)$ è l'incentro, % lang-it
\item $X(2)$ è il baricentro, % lang-it
\item $X(3)$ è il circocentro, % lang-it
\item $X(4)$ è l'ortocentro. % lang-it
\end{itemize}
Te cztery punkty opisujemy poniżej. % lang-pl
Questi quattro punti sono descritti qui di seguito. % lang-it

\begin{itemize}
\item $X(5)$ to środek okręgu dziewięciu punktów z faktu \ref{okrag_dziewieciu_punktow}, % lang-pl
\item $X(6)$ to punkt Lemoine'a (Grebego) z faktu \ref{punkt_lemoine}, % lang-pl
\item $X(7)$ to punkt Gergonne'a z~faktu \ref{punkt_gergonne}, % lang-pl
\item $X(8)$ to punkt Nagela z faktu \ref{punkt_nagela}, % lang-pl
\item $X(9)$ to mittenpunkt (punkt Lemoine'a trójkąta rozpiętego przez środki okręgów dopisanych; leży na prostych łączących środek ciężkości z punktem Gergonne'a, środek okręgu wpisanego z punktem Lemoine'a oraz ortocentrum ze środkiem Spiekera), % lang-pl
\item $X(10)$ to środek okręgu Spiekera z faktu \ref{punkt_spiekera}, % lang-pl
\item $X(11)$ to punkt Feuerbacha z twierdzenia \ref{punkt_feuerbacha}, % lang-pl
\item $X(13)$ to punkt Fermata z zadania \ref{punkt_fermata}, % lang-pl
\item $X(17)$, $X(18)$ to punkty Napoleona, % lang-pl
\item $X(39)$ to środek odcinka łączącego punkty Brocarda z definicji \ref{punkty_brocarda}. % lang-pl
\item $X(5)$ è il centro del cerchio dei nove punti del fatto \ref{okrag_dziewieciu_punktow}, % lang-it
\item $X(6)$ è il punto di Lemoine (Grebe) del fatto \ref{punkt_lemoine}, % lang-it
\item $X(7)$ è il punto di Gergonne del fatto \ref{punkt_gergonne}, % lang-it
\item $X(8)$ è il punto di Nagel del fatto \ref{punkt_nagela}, % lang-it
\item $X(9)$ è il mittenpunkt (il punto di Lemoine del triangolo formato dai centri degli excerchi; appartiene alle rette che uniscono il baricentro con il punto di Gergonne, l'incentro con il punto di Lemoine e l'ortocentro con il centro di Spieker), % lang-it
\item $X(10)$ è il centro del cerchio di Spieker del fatto \ref{punkt_spiekera}, % lang-it
\item $X(11)$ è il punto di Feuerbach del teorema \ref{punkt_feuerbacha}, % lang-it
\item $X(13)$ è il punto di Fermat dell'esercizio \ref{punkt_fermata}, % lang-it
\item $X(17)$, $X(18)$ sono i punti di Napoleone, % lang-it
\item $X(39)$ è il punto medio del segmento che unisce i punti di Brocard definiti in \ref{punkty_brocarda}. % lang-it
\end{itemize}

Lista jest długa, a jej końca nie widać. % lang-pl
La lista è lunga e sembra non avere fine. % lang-it

% eksperymentalne tłumaczenie na włoski

%

\subsection{Okrąg opisany i trzy symetralne} % lang-pl

\subsection{Circonferenza circoscritta e tre assi dei segmenti} % lang-it

\begin{proposition}
    \label{symetralne_przecinaja_sie}
    Trzy symetralne boków trójkąta przecinają się w jednym punkcie, środku okręgu opisanego na tym trójkącie. % lang-pl
    
    I tre assi dei lati di un triangolo si intersecano in un unico punto, il centro della circonferenza circoscritta al triangolo. % lang-it
\end{proposition}

Wynika to bezpośrednio z uwagi za definicją \ref{def_symetralna}: w trójkącie $\triangle ABC$, symetralna boku $AB$ (odpowiednio: $AC$) zawiera środki okręgów przechodzących przez punkty $A$ oraz $B$ (przez punkty $A$ oraz $C$), zatem trzecia symetralna musi przejść przez punkt wspólny dwóch pierwszych. % lang-pl
Ciò segue direttamente dall'osservazione successiva alla definizione \ref{def_symetralna}: nel triangolo $\triangle ABC$, l'asse del lato $AB$ (rispettivamente $AC$) contiene i centri delle circonferenze passanti per i punti $A$ e $B$ (rispettivamente per $A$ e $C$), dunque il terzo asse deve passare per il punto comune dei primi due. % lang-it


Trzy różne punkty współliniowe nie leżą na jednym okręgu, zaś dowolne trzy niewspółliniowe punkty leżą na dokładnie jednym okręgu. % lang-pl
 % Guzicki s. 15 % lang-pl
Hartshorne \cite[s. 16]{hartshorne2000} podaje to w formie ćwiczenia ze wskazówką, by spojrzeć na (IV.5). % lang-pl
Audin \cite[s. 61]{audin_2003} też, ale bez wskazówki. % lang-pl
Zetel \cite[s. 151]{zetel_2020} poda ciekawe uzasadnienie z trójkątami rzutowymi. % lang-pl
Inne odsyłacze to Neugebauer \cite[s. 19]{neugebauer_2018}. % lang-pl
Tre punti distinti e allineati non appartengono a una stessa circonferenza, mentre tre punti qualsiasi non allineati appartengono a un'unica circonferenza. % Guzicki s. 15 % lang-it


Hartshorne \cite[p.~16]{hartshorne2000} lo propone come esercizio, suggerendo di consultare (IV.5). % lang-it
Anche Audin \cite[p.~61]{audin_2003}, ma senza suggerimenti. % lang-it
Zetel \cite[p.~151]{zetel_2020} ne fornisce una dimostrazione interessante mediante triangoli proiettivi. % lang-it
Un altro riferimento è Neugebauer \cite[p.~19]{neugebauer_2018}. % lang-it
\begin{proposition}
	Środek okręgu opisanego na trójkącie leży wewnątrz tego trójkąta (odpowiednio: na jego brzegu, na zewnątrz trójkąta) wtedy i tylko wtedy, gdy trójkąt jest ostrokątny (odpowiednio: prostokątny, rozwartokątny). % lang-pl
	
	Il centro della circonferenza circoscritta a un triangolo si trova all'interno del triangolo (rispettivamente sul suo bordo, all'esterno del triangolo) se e solo se il triangolo è acutangolo (rispettivamente rettangolo, ottusangolo). % lang-it
\end{proposition}

Uogólnieniem faktu \ref{symetralne_przecinaja_sie} o współpękowości symetralnych boków trójkąta jest: % lang-pl

Una generalizzazione del fatto \ref{symetralne_przecinaja_sie} sulla concorrenza degli assi dei lati di un triangolo è il seguente risultato: % lang-it

\begin{theorem}[Carnota]
\index{twierdzenie!Carnota}%
\index{teorema!di Carnot}% % lang-it
\label{guzicki_6_13}%
	Dany jest trójkąt $ABC$ i punkty $D, E, F$ leżące odpowiednio na prostych $BC, CA, AB$. % lang-pl
	
	Niech prosta $k$ (odpowiednio: $l$, $m$) przechodzi przez punkt $D$ ($E$, $F$) i będzie prostopadła do prostej $BC$ ($CA$, $AB$). % lang-pl
	Siano dati un triangolo $ABC$ e punti $D$, $E$, $F$ appartenenti rispettivamente alle rette $BC$, $CA$, $AB$. % lang-it
	
	Wtedy proste $k$, $l$, $m$ mają punkt wspólny wtedy i tylko wtedy, gdy % lang-pl
	Sia la retta $k$ (rispettivamente $l$, $m$) passante per $D$ ($E$, $F$) e perpendicolare alla retta $BC$ ($CA$, $AB$). % lang-it
	
	Allora le rette $k$, $l$, $m$ sono concorrenti se e solo se % lang-it
	\begin{equation}
		|AF|^2 + |BD|^2 + |CE|^2 = |AE|^2 + |BF|^2 + |CD|^2.
	\end{equation}
\end{theorem}
% TODO: https://en.wikipedia.org/wiki/Carnot%27s_theorem_(perpendiculars)

Guzicki \cite[s. 176]{guzicki_2021} wyprowadza je z twierdzenia Pitagorasa, co pozwala mu dojść do wniosku, że okręgi opisany i wpisany oraz ortocentrum istnieją. % lang-pl

Guzicki \cite[p.~176]{guzicki_2021} lo deduce dal teorema di Pitagora, ottenendo così l'esistenza della circonferenza circoscritta, della circonferenza inscritta e dell'ortocentro. % lang-it
\index{twierdzenie!Pitagorasa}
\index{teorema!di Pitagora} % lang-it

%
% eksperymentalne tłumaczenie na włoski

%

\subsection{Okrąg wpisany i trzy dwusieczne} % lang-pl
\subsection{Cerchio inscritto e tre bisettrici} % lang-it

\begin{proposition}
    Trzy dwusieczne kątów trójkąta przecinają się w jednym punkcie, środku okręgu wpisanego w ten trójkąt. % lang-pl
    Le tre bisettrici degli angoli di un triangolo si incontrano in un unico punto, il centro del cerchio inscritto nel triangolo. % lang-it
\end{proposition}

Uzasadnienie jest kompletnie analogiczne do istnienia środku okręgu opisanego (tak zrobi Neugebauer \cite[s. 32]{neugebauer_2018}), ale można też podać fikuśny dowód dookoła: z twierdzenia o dwusiecznej i Cevy, jak Zetel \cite[s. 14]{zetel_2020}. % lang-pl
Hartshorne \cite[s. 16]{hartshorne_2000} podaje to w formie ćwiczenia ze wskazówką, by spojrzeć na (IV.4), z czym totalnie się zgadzamy. % lang-pl
La dimostrazione è del tutto analoga a quella dell'esistenza del centro del cerchio circoscritto (come fa Neugebauer \cite[p. 32]{neugebauer_2018}), ma si può anche fornire una dimostrazione più elegante: mediante il teorema della bisettrice e il teorema di Ceva, come in Zetel \cite[p. 14]{zetel_2020}. % lang-it
Hartshorne \cite[p. 16]{hartshorne_2000} lo presenta come esercizio con il suggerimento di guardare (IV.4), e siamo pienamente d'accordo. % lang-it

Po angielsku mamy krótkie określenia \emph{excenter, excircle, incenter, incircle}. % lang-pl
In inglese si usano i termini brevi \emph{excenter, excircle, incenter, incircle}. % lang-it

% TODO: długość dwusiecznej, Zetel s. 32

\begin{proposition}
    Niech $r$ i $R$ będą promieniami okręgów wpisanego i opisanego na trójkącie o kątach $\alpha, \beta, \gamma$. % lang-pl
    Wtedy % lang-pl
    Siano $r$ e $R$ i raggi del cerchio inscritto e circoscritto di un triangolo con angoli $\alpha, \beta, \gamma$. % lang-it
    Allora % lang-it
    \begin{equation}
        \frac r R = 4 \sin \frac \alpha 2 \sin \frac \beta 2 \sin \frac \gamma 2.
    \end{equation}
\end{proposition}

Trygonometryczne rozwiązanie opisze Guzicki \cite[s. 351]{guzicki_2021}. % lang-pl
Una dimostrazione trigonometrica è presentata da Guzicki \cite[p. 351]{guzicki_2021}. % lang-it

%
% eksperymentalne tłumaczenie na włoski

%

\subsection{Okręgi dopisane} % lang-pl
\index{okrąg dopisany}% % lang-pl
\subsection{Circonferenze exinscritte} % lang-it
\index{circonferenza exinscritta}% % lang-it

Oprócz okręgu wpisanego w trójkąt są jeszcze trzy inne okręgi styczne do wszystkich trzech prostych, na których leżą boki trójkąta. % lang-pl
Nazywamy je okręgami dopisanymi (\emph{excircle} albo \emph{escribed circle} po angielsku), pisze o nich Audin \cite[s. 98]{audin_2003}. % lang-pl
Oltre alla circonferenza inscritta in un triangolo esistono altre tre circonferenze tangenti a tutte e tre le rette contenenti i lati del triangolo. % lang-it
Sono dette circonferenze exinscritte (\emph{excircle} oppure \emph{escribed circle} in inglese); Audin ne parla a \cite[p. 98]{audin_2003}. % lang-it

Środek okręgu wpisanego znajduje się na przecięciu dwusiecznych wewnętrznych kątów trójkąta. % lang-pl
Analogicznie środki okręgów dopisanych znajdują się na przecięciu dwusiecznej kąta wewnętrznego przy ustalonym wierzchołku i dwusiecznych kątów zewnętrznych przy pozostałych wierzchołkach. % lang-pl
Il centro della circonferenza inscritta è il punto d'incontro delle bisettrici interne del triangolo. % lang-it
Analogamente, i centri delle circonferenze exinscritte si trovano all'intersezione della bisettrice dell'angolo interno in un vertice fissato e delle bisettrici degli angoli esterni negli altri due vertici. % lang-it
% TODO: https://upload.wikimedia.org/wikipedia/commons/thumb/6/6f/Incircle_and_Excircles.svg/960px-Incircle_and_Excircles.svg.png
Ich promienie można łatwo wyznaczyć przez porównanie pól trójkątów, co prowadzi do wzoru: % lang-pl
I loro raggi si determinano facilmente confrontando aree di triangoli, ottenendo la formula % lang-it
\begin{equation}
	r_a = \frac{S_{ABC}}{p - a},
\end{equation}
oraz dwóch kolejnych dla $r_b$ i $r_c$, jak pokazuje Zetel \cite[s. 18]{zetel_2020}. % lang-pl
Tutaj $p$ oznacza tak jak zawsze połowę obwodu. % lang-pl
Czynnik $p - a$ sugeruje związek ze wzorem Herona i tak istotnie jest, mamy % lang-pl
e le analoghe formule per $r_b$ e $r_c$, come mostra Zetel \cite[p. 18]{zetel_2020}. % lang-it
Qui $p$ indica, come di consueto, il semiperimetro. % lang-it
Il fattore $p-a$ suggerisce un legame con la formula di Erone, e infatti % lang-it
\begin{equation}
	S_{ABC}^2 = r \cdot r_a \cdot r_b \cdot r_c.
\end{equation}

Zależności między czterema promieniami jest bez liku, mamy też: % lang-pl
Le relazioni tra questi quattro raggi sono numerosissime; ad esempio % lang-it
\begin{equation}
	\frac{1}{r} = \frac {1} {r_a} + \frac {1} {r_b} + \frac {1} {r_c}.
\end{equation}

% TODO: Wzory na promienie okręgów wpisanych, dopisanych.
\todofoot{$4R = r_a + r_b + r_c - r$} % Coxeter, s. 13; dopisz też bend okręgi Kartezjusza
% eksperymentalne tłumaczenie na włoski

%

\subsection{Ortocentrum i trzy wysokości}

\begin{definition}[wysokość]
\index{wysokość trójkąta}%
    \pol{Niech $\triangle ABC$ będzie trójkątem, w którym przez punkt $C$ poprowadzono prostą prostopadłą do boku $AB$.}
    \ita{Sia $\triangle ABC$ un triangolo e si tracci per il punto $C$ la retta perpendicolare al lato $AB$.}
    \pol{Odcinek leżący na tej prostej, którego jeden koniec znajduje się w punkcie $C$, zaś drugi leży na boku $AB$, nazywamy wysokością opuszczoną z punktu $C$, drugi jego koniec to spodek wysokości.}
    \ita{Il segmento su tale retta avente un estremo nel punto $C$ e l'altro sul lato $AB$ si chiama altezza relativa al punto $C$, e il secondo estremo è il piede dell'altezza.}
\end{definition}

\pol{Po angielsku mamy \emph{altitude} oraz \emph{foot of the altitude}.}
\ita{In inglese si usano i termini \emph{altitude} e \emph{foot of the altitude}.}

% TODO: Sometimes an arbitrary edge is chosen to be the base, in which case the opposite vertex is called the apex; the shortest segment between the base and apex is the height. The area of a triangle equals one-half the product of height and base length.

\begin{proposition}
\label{wysokosci_przecinaja_sie}%
	\pol{Trzy wysokości trójkąta przecinają się w jednym punkcie zwanym ortocentrum.}
	\ita{Le tre altezze di un triangolo si intersecano in un unico punto chiamato ortocentro.}
\index{ortocentrum}%
\end{proposition}

\pol{Wysokości trójkąta rozwartokątnego nie przecinają się, musimy je więc zastąpić przez proste, na których leżą.}
\ita{Le altezze di un triangolo ottusangolo non si intersecano all'interno del triangolo, quindi si considerano le rette su cui giacciono.}

\pol{Dowodów we współczesnej literaturze nie brakuje: są u Pompego \cite[s. 38]{pompe_2022}, gdzie zmyślnie używa równoległoboków, Guzickiego \cite[s. 218]{guzicki_2021}, Zetela \cite[s. 14, 151]{zetel_2020}.}
\ita{Le dimostrazioni nella letteratura moderna non mancano: si trovano in Pompe \cite[p.~38]{pompe_2022}, dove si usano in modo ingegnoso i parallelogrammi, in Guzicki \cite[p.~218]{guzicki_2021}, e in Zetel \cite[pp.~14, 151]{zetel_2020}.}
\pol{Warto też przeczytać tekst Hartshorne'a \cite[s. 54]{hartshorne2000} oraz Audina \cite[s. 61]{audin_2003}.}
\ita{Vale anche la pena leggere Hartshorne \cite[p.~54]{hartshorne2000} e Audin \cite[p.~61]{audin_2003}.}

\pol{Ortocentrum (\emph{orthocenter}) będzie czasami nazywane środkiem ortycznym, ale nazwa ta nie utrzyma się długo.}
\ita{L'ortocentro (\emph{orthocenter}) viene talvolta chiamato centro ortico, ma questo termine non si affermerà a lungo.}

\begin{proposition}
    \pol{Dany jest trójkąt $ABC$ oraz jego ortocentrum $H$.}
    \ita{Sia dato un triangolo $ABC$ e il suo ortocentro $H$.}
    \pol{Odbicia symetryczne punktu $H$ względem prostych zawierających boki trójkąta $ABC$ leżą na okręgu opisanym na tym trójkącie.}
    \ita{Le simmetrie del punto $H$ rispetto alle rette contenenti i lati del triangolo $ABC$ giacciono sulla circonferenza circoscritta al triangolo.}
\end{proposition}

\pol{Dowód tego prostego stwierdzenia przedstawia Neugebauer \cite[s. 28]{neugebauer_2018}.}
\ita{La dimostrazione di questo semplice fatto è presentata in Neugebauer \cite[p.~28]{neugebauer_2018}.}

%
% eksperymentalne tłumaczenie na włoski

%

\subsection{Środek ciężkości i trzy środkowe} % lang-pl
\subsection{Baricentro e tre mediane} % lang-it
\begin{definition}[mediana] % lang-it
\index{mediana}% % lang-it
    Dato un triangolo. % lang-it
    Il segmento che unisce uno qualsiasi dei suoi vertici con il punto medio del lato opposto si chiama mediana. % lang-it
\end{definition} % lang-it


\begin{definition}[środkowa] % lang-pl
\index{środkowa}% % lang-pl
    Dany jest pewien trójkąt. % lang-pl
    Odcinek łaczący którykolwiek z jego wierzchołków ze środkiem przeciwległego boku nazywamy środkową. % lang-pl
\end{definition} % lang-pl
Da non confondere con il segmento medio! % lang-it

Nie mylić z linią środkową! % lang-pl
Środkowe dzielą trójkąt na sześć mniejszych o równych polach. % lang-pl
Le mediane dividono il triangolo in sei triangoli più piccoli di uguale area. % lang-it
% TODO: rysunek równych sześciu pól.

% TODO: długość środkowej, Zetel s. 32

\begin{proposition}
    Niech $m_a, m_b, m_c$ oznacza długości trzech środkowych trójkąta, zaś $p$ połowę jego obwodu. % lang-pl
    
    Wtedy % lang-pl
    Siano $m_a, m_b, m_c$ le lunghezze delle tre mediane di un triangolo e $p$ il semiperimetro. % lang-it
    
    Allora % lang-it
    \begin{equation}
        \frac 32 p \le m_a + m_b + m_c \le 2p.
    \end{equation}
    % Posamentier, Alfred S., and Salkind, Charles T., Challenging Problems in Geometry, Dover, 1996: pp. 86–87.
\end{proposition}

Napisze o tym Coxeter \cite[s. 26]{coxeter_1967}. % lang-pl

Guzicki \cite[s. 357]{guzicki_2021} wytłumaczy dodatkowo, czemu żadnej ze stałych: $\frac 34$, $1$ nie można poprawić. % lang-pl
Ne scriverà Coxeter \cite[p. 26]{coxeter_1967}. % lang-it
Guzicki \cite[p. 357]{guzicki_2021} spiegherà inoltre perché nessuna delle costanti $\frac 34$, $1$ può essere migliorata. % lang-it


Mamy też podobną nierówność: % lang-pl
Abbiamo anche una disuguaglianza simile: % lang-it

\begin{proposition}
    Niech $a, b, c$ oznacza długości boków trójkąta. % lang-pl
    
    Wtedy % lang-pl
    Siano $a$, $b$, $c$ le lunghezze dei lati di un triangolo. % lang-it
    
    Allora % lang-it
    \begin{equation}
        \frac 32 \le \frac{a}{b+c} + \frac{b}{a+c} + \frac{c}{a+b} < 2.
    \end{equation}
\end{proposition}

Lewa połowa jest prawdziwa dla wszystkich dodatnich liczb $a, b, c$ i nosi nazwę nierówności (Alfreda Mortimera) Nesbitta, który znajdzie ją około 1902 roku. % lang-pl

La disuguaglianza a sinistra vale per tutti i numeri positivi $a, b, c$ ed è nota come disuguaglianza di (Alfred Mortimer) Nesbitt, che la troverà intorno al 1902. % lang-it

\begin{proposition}
\label{srodkowe_przecinaja_sie}%
\index{środek ciężkości}% % lang-pl
    Trzy środkowe trójkąta przecinają się w jednym punkcie zwanym środkiem ciężkości i dzielą się w~stosunku $2$ do $1$, licząc od wierzchołków. % lang-pl
\index{baricentro}% % lang-it
    
    Le tre mediane di un triangolo si incontrano in un unico punto, chiamato baricentro, e vengono divise da esso nel rapporto $2$ a $1$, contando dai vertici. % lang-it
\end{proposition}
% % Coxeter, Introduction to Geometry, s. 10 <- przeczytaj to, nie tylko cytuj! + ćwiczenia: 3/4 <= 1
% Neugebauer s. 54: oznaczenie G, bo gravis - ciężki
% Zetel, s. 25 z tw. Cevy

Środkowe to po angielsku \emph{medians}, przecinają się w \emph{centroid}. % lang-pl

Polska nazwa weźmie się z tego, że środek ciężkości fizycznego modelu trójkąta wykonanego z jednolitego materiału znajduje się właśnie tam. % lang-pl
Po angielsku środkowe nazywają się \emph{medians}, a ich punkt przecięcia \emph{centroid}. % lang-it

Angielska nazwa powstanie w 1814 roku z powodu niechęci do tego, co nie jest czysto geometryczne; niechęci, której nie będzie widać w innych europejskich językach. % lang-pl
Il nome polacco deriva dal fatto che il centro di gravità di un modello fisico di triangolo realizzato con materiale omogeneo si trova proprio in quel punto. % lang-it
Il termine inglese nascerà nel 1814 a causa di una certa avversione verso ciò che non è puramente geometrico; un'avversione che non si osserverà nelle altre lingue europee. % lang-it


(Środek ciężkości brzegu trójkąta nazywa się punktem Spiekera, patrz fakt \ref{punkt_spiekera}). % lang-pl
(Il centro di gravità del bordo di un triangolo si chiama punto di Spieker, vedi il fatto \ref{punkt_spiekera}). % lang-it
\index{punkt!Spiekera}%

Hartshorne \cite[s. 52-54]{hartshorne2000} wnioskuje \ref{srodkowe_przecinaja_sie} z faktu \ref{hartshorne_52} (później zaś \cite[s. 119-120]{hartshorne2000} powtórzy dowód z~maszynerią geometrii analitycznej). % lang-pl

Hartshorne \cite[p. 52-54]{hartshorne2000} deduce \ref{srodkowe_przecinaja_sie} dal fatto \ref{hartshorne_52} (più tardi \cite[p. 119-120]{hartshorne2000} ripeterà la dimostrazione usando gli strumenti della geometria analitica). % lang-it
\index[persons]{Archimedes}% % lang-pl
\index[persons]{Archimede}% % lang-it
Podobnie zrobią Pompe \cite[s. 41]{pompe_2022}, Guzicki \cite[s. 220]{guzicki_2021} albo Bogdańska, Neugebauer \cite[s. 53]{guzicki_2021}. % lang-pl

Zetel \cite[s. 14, 25]{zetel_2020} pokaże, że twierdzenie Cevy daje natychmiast dowód istnienia środka ciężkości. % lang-pl
Faranno lo stesso Pompe \cite[p. 41]{pompe_2022}, Guzicki \cite[p. 220]{guzicki_2021} e Bogdańska, Neugebauer \cite[p. 53]{guzicki_2021}. % lang-it

\index{twierdzenie!Cevy}% % lang-pl
Zetel \cite[p. 14, 25]{zetel_2020} mostrerà che il teorema di Ceva fornisce immediatamente una dimostrazione dell'esistenza del baricentro. % lang-it
\index{teorema!di Ceva}% % lang-it
% TODO: Ich długość można obliczyć z https://en.wikipedia.org/wiki/Apollonius%27s_theorem => to jest z https://en.wikipedia.org/wiki/Median_(geometry)#Formulas_involving_the_medians'_lengths
Coxeter \cite[s. 26]{coxeter_1967} pokaże, jak otrzymać ten sam wynik z (VI.2) oraz (VI.4). % lang-pl
Coxeter \cite[p. 26]{coxeter_1967} mostrerà come ottenere lo stesso risultato da (VI.2) e (VI.4). % lang-it


Historia faktów \ref{wysokosci_przecinaja_sie} oraz \ref{srodkowe_przecinaja_sie} ma wiele elementów wspólnych, więc omówimy je razem. % lang-pl
La storia dei fatti \ref{wysokosci_przecinaja_sie} e \ref{srodkowe_przecinaja_sie} presenta molti elementi in comune, perciò li discuteremo insieme. % lang-it

Żaden z nich nie pojawi się w Elementach Euklidesach i właściwie nie wiadomo, kto odkryje je pierwszy. % lang-pl
Nessuno dei due compare negli Elementi di Euclide e non si sa realmente chi li abbia scoperti per primo. % lang-it

Uczeni arabscy stwierdzą, że był to Archimedes, choć nie ma ku temu niezbitych dowodów. % lang-pl
Gli studiosi arabi affermeranno che fu Archimede, anche se non esistono prove definitive. % lang-it
\index[persons]{Archimede}% % lang-it

\index[persons]{Archimedes}% % lang-pl
Pappus będzie świadom istnienia ortocentrum, ale najstarsze znane uzasadnienie (w komentarzu al-Nasawiego) przyjdzie na świat dopiero w XI wieku. % lang-pl
Pappo conoscerà l'esistenza dell'ortocentro, ma la più antica dimostrazione nota (in un commento di al-Nasawi) apparirà soltanto nell'XI secolo. % lang-it
\index[persons]{Pappus}% % lang-pl
\index[persons]{Pappo}% % lang-it
\index[persons]{al-Nasawi, Ali ibn Ahmad}%

Fakt \ref{srodkowe_przecinaja_sie} ma trójwymiarową kontynuację: % lang-pl
Il fatto \ref{srodkowe_przecinaja_sie} possiede un'analoga versione tridimensionale: % lang-it
\begin{proposition} % lang-it
    I quattro segmenti che uniscono i vertici di un tetraedro con i baricentri delle facce opposte si incontrano in un unico punto, che li divide nel rapporto $3$ a $1$. % lang-it
\end{proposition} % lang-it


\begin{proposition} % lang-pl
    Cztery odcinki łączące wierzchołki czworościanu ze środkami przeciwległych ścian przecinają się w~jednym punkcie, który dzieli je w stosunku $3$ do $1$. % lang-pl
\end{proposition} % lang-pl
Alcuni usano l'espressione ``teorema di Commandino'', poiché Federico Commandino scriverà nel 1565 l'opera \emph{De centro gravitatis solidorum} (Sui centri di gravità dei solidi), pur non essendo necessariamente il primo ad averlo scoperto. % lang-it
\index{teorema!di Commandino}% % lang-it
\index[persons]{Commandino, Federico}% % lang-it

Niektórzy stosują określenie ,,twierdzenie Commandino'', ponieważ Federico Commandino napisze w 1565 roku pracę \emph{De centro gravitatis solidorum} (O środkach ciężkości brył), chociaż może nie być pierwszym, który je odkrył. % lang-pl
Si sospetta che Francesco Maurolico lo conoscesse già in precedenza. % lang-it

\index{twierdzenie!Commandino}% % lang-pl
\index[persons]{Commandino, Federico}% % lang-pl
Podejrzewa się, że Francesco Maurolico pozna je wcześniej. % lang-pl
(Friederich Eduard Reusch znajdzie uogólnienie, które w zdegenerowanym przypadku prowadzi znowu do twierdzenia \ref{theorem_varignon} Varignona). % lang-pl
(Friederich Eduard Reusch troverà una generalizzazione che, in un caso degenere, conduce nuovamente al teorema di Varignon \ref{theorem_varignon}). % lang-it
\index[persons]{Reusch, Friederich Eduard}%

%
% eksperymentalne tłumaczenie na włoski

%

\subsection{Prosta Wallace'a-Simsona} % lang-pl
\subsection{Retta di Wallace-Simson} % lang-it

\begin{proposition}
[twierdzenie Wallace'a] % lang-pl
[teorema di Wallace] % lang-it
	Dany jest trójkąt $ABC$ oraz punkt $P$. % lang-pl
	Niech $P_a$, $P_b$, $P_c$ oznacza rzuty punktu $P$ na proste $BC$, $AC$, $AB$. % lang-pl
	Wtedy następujące warunki są równoważne: punkt $P$ leży na okręgu opisanym na $\triangle ABC$; punkty $P_a$, $P_b$, $P_c$ są współliniowe. % lang-pl
	Siano dati un triangolo $ABC$ e un punto $P$. % lang-it
	Siano $P_a$, $P_b$, $P_c$ le proiezioni del punto $P$ sulle rette $BC$, $AC$, $AB$. % lang-it
	Allora le seguenti condizioni sono equivalenti: il punto $P$ appartiene alla circonferenza circoscritta a $\triangle ABC$; i punti $P_a$, $P_b$, $P_c$ sono allineati. % lang-it
\end{proposition}

Kiedy punkt $P$ leży na wspomnianym okręgu, prostą przez punkty $P_a$, $P_b$ oraz $P_c$ nazywa się prostą Simsona; na cześć Roberta Simsona, szkockiego matematyka żyjącego w latach 1687-1768; chociaż pierwszą pracę na jej temat opublikuje Wallace w 1799 roku. % lang-pl
Quando il punto $P$ appartiene alla circonferenza suddetta, la retta passante per i punti $P_a$, $P_b$ e $P_c$ è detta retta di Simson, in onore di Robert Simson, matematico scozzese vissuto tra il 1687 e il 1768, sebbene il primo lavoro sull'argomento sia stato pubblicato da Wallace nel 1799. % lang-it
\index[persons]{Wallace, ?}

Być może właściwsze byłoby pisać ,,prosta Wallace'a-Simsona''. % lang-pl
Forse sarebbe più corretto parlare di ``retta di Wallace-Simson''. % lang-it

Współliniowość spodków punktu $P$ na boki trójkąta jest twierdzeniem z dowodem u Zetela \cite[s. 55]{zetel_2020}; ćwiczeniem u Audina \cite[s. 104]{audin_2003}, Hartshorne'a \cite[s. 61]{hartshorne2000} i Neugebauera \cite[s. 25]{neugebauer_2018}. % lang-pl
L'allineamento dei piedi delle perpendicolari dal punto $P$ ai lati del triangolo è presentato come teorema con dimostrazione da Zetel \cite[p. 55]{zetel_2020}; come esercizio da Audin \cite[p. 104]{audin_2003}, Hartshorne \cite[p. 61]{hartshorne2000} e Neugebauer \cite[p. 25]{neugebauer_2018}. % lang-it
% \todofoot{HARTSHORNE PRZECHODZI DO PUNKTU MIQUELA}
Znajdzie się też w \cite[s. 24]{komisarczyk_2024}. % lang-pl
Pojawi się wielokrotnie w Delcie: w artykułach Michała Kiezy ($\Delta_{11}^9$), Joanny Jaszuńskiej ($\Delta_{15}^{10}$), Dominika Burka i Tomasza Cieśli ($\Delta_{16}^{11}$, z twierdzeniem Steinera). % lang-pl
Compare anche in \cite[p. 24]{komisarczyk_2024}. % lang-it
Compare più volte anche su \emph{Delta}: negli articoli di Michał Kieza ($\Delta_{11}^9$), Joanna Jaszuńska ($\Delta_{15}^{10}$), Dominik Burek e Tomasz Cieśla ($\Delta_{16}^{11}$, con il teorema di Steiner). % lang-it

\begin{proposition}
[twierdzenie Steinera] % lang-pl
[teorema di Steiner] % lang-it
	Dany jest trójkąt $ABC$ oraz punkt $P$ leżący na okręgu opisanym na trójkącie. % lang-pl
	Wtedy obrazy punktu $P$ w symetrii względem boków trójkąta leżą na jednej prostej, która przechodzi przez ortocentrum trójkąta. % lang-pl
	Siano dati un triangolo $ABC$ e un punto $P$ appartenente alla sua circonferenza circoscritta. % lang-it
	Allora le immagini del punto $P$ rispetto alle simmetrie aventi come assi i lati del triangolo giacciono su una stessa retta, che passa per l'ortocentro del triangolo. % lang-it
\end{proposition}

% PROSTA STEINERA? Audin ta sama strona
% prosta Steinera -> Znajdzie się też w \cite[s. 24]{komisarczyk_2024}.
% http://users.math.uoc.gr/~pamfilos/eGallery/problems/SteinerLine.html

Z twierdzenia o prostej Wallace'a-Simsona wynika, jak u Zetela \cite[s. 57]{zetel_2020}: % lang-pl
Dal teorema della retta di Wallace-Simson segue, come in Zetel \cite[p. 57]{zetel_2020}: % lang-it

\begin{proposition}
[twierdzenie Salmona] % lang-pl
[teorema di Salmon] % lang-it
\index{twierdzenie!Salmona}% % lang-pl
	Dany jest okrąg oraz trzy jego różne cięciwy $PA$, $PB$, $PC$ takie, że przekrojem okręgów na średnicach $PA$, $PB$ (odpowiednio: $PB$, $PC$ i $PA$, $PC$) są punkty $P$, $M$ (odpowiednio: $P$, $K$ oraz $P$, $L$). % lang-pl
	Wtedy punkty $K$, $L$, $M$ są współliniowe. % lang-pl
\index{teorema!di Salmon}% % lang-it
	Sono dati un cerchio e tre sue corde distinte $PA$, $PB$, $PC$ tali che l'intersezione dei cerchi aventi per diametri $PA$, $PB$ (rispettivamente $PB$, $PC$ e $PA$, $PC$) sia costituita dai punti $P$, $M$ (rispettivamente $P$, $K$ e $P$, $L$). % lang-it
	Allora i punti $K$, $L$, $M$ sono allineati. % lang-it
\end{proposition}

%

%